%! TeX root = TeoríaDeCuerdas.tex

\bgimg{img/VibratingStringsBackground.png}
\sectiontitle{SOBRE LA TEORÍA DE CUERDAS}

\begin{frame}{
%%%%Title%%%%
    Sobre la teoría de cuerdas
}
%%%%Content%%%%
  \onslide*<1>{
    La \textbf{teoría de cuerdas} propone esencialmente que los componentes fundamentales del
    universo no son partículas puntuales, sino \textbf{filamentos de energía} que definen las
    características de los elementos subatómicos, y sus interacciones.
    Basados en la frecuencia y forma de sus oscilaciones.
  }

  \onslide*<3>{
    Esta teoría es capaz, matemáticamente de...
    \begin{itemize}
      \item Predecir las diferencias de masa entre partículas fundamentales.
      \item Describir la interacción gravitatoria.
      \item Proveer supersimetría entre partículas bosónicas y fermiónicas.
      \item Predecir el comportamiento relativista del universo
    \end{itemize}
  }
  
  \onslide*<2>{
    Desde su origen hasta la fecha ha pasado por diversas revisiones, refinándola
    y conceptualizando elementos más allá de las propias cuerdas. Esto llevó a los
    diferentes planteamientos que, al final, se demostraron ser diferentes perspectivas
    de una misma teoría.
  }

  \onslide<1-3>{
    \vfill
    \begin{center}
      \begin{adjustbox}{max width=0.7\textwidth}
        \input{img/VibratingStrings}
      \end{adjustbox}
    \end{center}
  }

\end{frame}
